\section{Introduction}

\label{sec:intro}
% ============================================================

Cross-chain bridges must hold custody of assets on source chains while issuing
representations on destination chains. The security of billions of dollars in locked
value depends on how the custodial key is managed. Three approaches dominate
production deployments:

\begin{enumerate}[nosep]
  \item \textbf{Multisig}: $k$-of-$n$ on-chain multisignature. Transparent but
    expensive (linear gas cost in $k$) and chain-specific (not all chains support
    native multisig).
  \item \textbf{Trusted committee}: A centralized signer or small committee holds
    the private key. Cheap verification but catastrophic failure mode.
  \item \textbf{Threshold signatures}: $t$-of-$n$ MPC-based signing produces a
    single standard signature indistinguishable from a single-signer output.
    Constant verification cost. No on-chain multisig infrastructure required.
\end{enumerate}

Threshold signatures are the only approach that achieves both constant on-chain cost
and distributed trust. However, the blockchain ecosystem uses two incompatible
signature families: ECDSA (secp256k1) for Ethereum, Bitcoin legacy, TRON, and XRPL;
and Schnorr/EdDSA (Ed25519 or BIP-340) for Solana, Cardano, TON, Polkadot, Cosmos,
and Bitcoin Taproot. No single threshold protocol covers both.

\paragraph{Contribution.} We present a dual-protocol threshold framework:

\begin{itemize}[nosep]
  \item \textbf{FROST}~\cite{komlo2020frost} for Schnorr/EdDSA chains: 2-round
    threshold signing, 64-byte signatures, compatible with Ed25519 and BIP-340.
  \item \textbf{CGGMP21}~\cite{canetti2021cggmp} for ECDSA chains: threshold ECDSA
    with identifiable abort, 65-byte signatures, compatible with secp256k1.
  \item \textbf{Shared infrastructure}: common DKG, key refresh, signer rotation,
    and monitoring layer for both protocols.
  \item \textbf{Formal composition}: proof that sharing infrastructure between the
    two protocols does not degrade security of either.
\end{itemize}

The framework is deployed in Lux Teleport, securing cross-chain transfers across
18 native bridge programs and 270 registered chain IDs. On-chain verification uses
a single MPC group address---the aggregate public key from DKG---rather than
individual signer keys: $\texttt{ecrecover}(\text{digest}, \sigma) = \texttt{mpcGroupAddress}$.
Individual signer identities exist for accountability and identifiable abort only.

\subsection{Notation}

We use the following notation throughout. Let $\mathbb{G}$ denote a cyclic group of
prime order $q$ with generator $G$. For a set $S$ of signers with $|S| = t$, the
Lagrange coefficient of party $i$ in $S$ is $\lambda_{i,S} = \prod_{j \in S, j \neq i}
\frac{j}{j - i} \bmod q$. We write $H : \{0,1\}^* \to \mathbb{F}_q$ for a
collision-resistant hash function modeled as a random oracle. $[n]$ denotes
$\{1, \ldots, n\}$. $x \xleftarrow{\$} S$ denotes uniform sampling from set $S$.

% ============================================================
